% GENERATED FILE. Edit canonical lessons, then run npm run generate:books.
% canonical-source-hash: fnv1a64:ab918cfca5fb2798
% canonical-lessons: PT-C483-espirrar, PT-C483-espirro, PT-C483-alergico, PT-C483-comichao, PT-C483-cocar

\chapter{Espirros e comichão}
\label{ch:pt-espirros-e-comichao}
\hlchaptermodality{483}

\begin{chapteropening}
\textbf{By the end of this chapter:} \emph{I can describe sneezing, itching and allergies.}

Complete the last lesson of chapter 483: I can describe sneezing, itching and allergies.
\end{chapteropening}

\section[espirrar]{espirrar — to sneeze}

\label{lesson:PT-C483-espirrar}



\begin{quote}
Before the new one: say the Portuguese for heel, then the Portuguese for navel.
\end{quote}

\subsection*{You'll want to know: espirrar}
\textbf{espirrar} — \textquotedblleft{}to sneeze\textquotedblright{}.

\begin{quote}
Na primavera, o meu pai espirra o dia todo.
\end{quote}

\begin{quote}
\emph{In spring, my father sneezes all day.}
\end{quote}

Past: \emph{espirrei}, \emph{espirrou}.

\begin{cousinweb}
\emph{espirrar} is inherited from Latin \emph{exspīrāre} \textquotedblleft{}to breathe out\textquotedblright{} (EtymDB), built on \emph{spīrāre} \textquotedblleft{}to breathe\textquotedblright{}.
\end{cousinweb}

\begin{grammarlens}[title={What you've built}]
One more word for telling, explaining and describing at B1.
\end{grammarlens}

\subsection*{Your turn}
\begin{itemize}
\raggedright
  \item \emph{Say it:} \emph{espirrar}
  \item \emph{Say it:} \emph{espirrar}, once more
  \item \emph{Say it:} a sentence of your own with \emph{espirrar}
  \item \emph{Recall:} say the Portuguese for heel, then the Portuguese for navel, then say \emph{espirrar} again
\end{itemize}

\begin{tcolorbox}[breakable,colback=teal!4,colframe=teal!35!black,title={Before you move on}]
What does espirrar mean? (\textquotedblleft{}To sneeze\textquotedblright{}.) Say it once more.
\end{tcolorbox}
\section[o espirro]{o espirro — sneeze}

\label{lesson:PT-C483-espirro}



\begin{quote}
Before the new one: say the Portuguese for navel, then the Portuguese for to sneeze.
\end{quote}

\subsection*{You'll want to know: o espirro}
\textbf{o espirro} — \textquotedblleft{}sneeze\textquotedblright{}.

\begin{quote}
Ouvi um espirro na sala ao lado.
\end{quote}

\begin{quote}
\emph{I heard a sneeze in the next room.}
\end{quote}

After a sneeze, people in Portugal say \emph{santinho!}

\begin{cousinweb}
\emph{espirro} comes from \emph{espirrar} \textquotedblleft{}to sneeze\textquotedblright{}, inherited from Latin \emph{exspīrāre} \textquotedblleft{}to breathe out\textquotedblright{} (EtymDB).
\end{cousinweb}

\begin{grammarlens}[title={What you've built}]
One more word for telling, explaining and describing at B1.
\end{grammarlens}

\subsection*{Your turn}
\begin{itemize}
\raggedright
  \item \emph{Say it:} \emph{o espirro}
  \item \emph{Say it:} \emph{o espirro}, once more
  \item \emph{Say it:} a sentence of your own with \emph{o espirro}
  \item \emph{Recall:} say the Portuguese for navel, then the Portuguese for to sneeze, then say \emph{o espirro} again
\end{itemize}

\begin{tcolorbox}[breakable,colback=teal!4,colframe=teal!35!black,title={Before you move on}]
What does o espirro mean? (\textquotedblleft{}Sneeze\textquotedblright{}.) Say it once more.
\end{tcolorbox}
\section[alérgico]{alérgico — allergic}

\label{lesson:PT-C483-alergico}



\begin{quote}
Before the new one: say the Portuguese for to sneeze, then the Portuguese for sneeze.
\end{quote}

\subsection*{You'll want to know: alérgico}
\textbf{alérgico} — \textquotedblleft{}allergic\textquotedblright{}.

\begin{quote}
Sou alérgico a gatos.
\end{quote}

\begin{quote}
\emph{I am allergic to cats.}
\end{quote}

\emph{Ser alérgico a} — to be allergic to. The illness is \emph{a alergia}.

\begin{cousinweb}
\emph{alérgico} is built on \emph{alergia}, which is usually traced to German \emph{Allergie}, a medical coinage from Greek \emph{állos} \textquotedblleft{}other\textquotedblright{} and \emph{érgon} \textquotedblleft{}work\textquotedblright{}.
\end{cousinweb}

\begin{grammarlens}[title={What you've built}]
One more word for telling, explaining and describing at B1.
\end{grammarlens}

\subsection*{Your turn}
\begin{itemize}
\raggedright
  \item \emph{Say it:} \emph{alérgico}
  \item \emph{Say it:} \emph{alérgico}, once more
  \item \emph{Say it:} a sentence of your own with \emph{alérgico}
  \item \emph{Recall:} say the Portuguese for to sneeze, then the Portuguese for sneeze, then say \emph{alérgico} again
\end{itemize}

\begin{tcolorbox}[breakable,colback=teal!4,colframe=teal!35!black,title={Before you move on}]
What does alérgico mean? (\textquotedblleft{}Allergic\textquotedblright{}.) Say it once more.
\end{tcolorbox}
\section[a comichão]{a comichão — itch}

\label{lesson:PT-C483-comichao}



\begin{quote}
Before the new one: say the Portuguese for sneeze, then the Portuguese for allergic.
\end{quote}

\subsection*{You'll want to know: a comichão}
\textbf{a comichão} — \textquotedblleft{}itch\textquotedblright{}.

\begin{quote}
A picada do mosquito deu-me comichão.
\end{quote}

\begin{quote}
\emph{The mosquito bite made me itch.}
\end{quote}

In Brazil: \emph{a coceira}. \emph{Ter comichão} — to itch.

\begin{cousinweb}
\emph{comichão} is usually traced to Latin \emph{comēstiō} \textquotedblleft{}a consuming\textquotedblright{}, from \emph{comedere} \textquotedblleft{}to eat up\textquotedblright{} (Lewis \& Short): an itch feels as if the skin were being eaten.
\end{cousinweb}

\begin{grammarlens}[title={What you've built}]
One more word for telling, explaining and describing at B1.
\end{grammarlens}

\subsection*{Your turn}
\begin{itemize}
\raggedright
  \item \emph{Say it:} \emph{a comichão}
  \item \emph{Say it:} \emph{a comichão}, once more
  \item \emph{Say it:} a sentence of your own with \emph{a comichão}
  \item \emph{Recall:} say the Portuguese for sneeze, then the Portuguese for allergic, then say \emph{a comichão} again
\end{itemize}

\begin{tcolorbox}[breakable,colback=teal!4,colframe=teal!35!black,title={Before you move on}]
What does a comichão mean? (\textquotedblleft{}Itch\textquotedblright{}.) Say it once more.
\end{tcolorbox}
\section[coçar]{coçar — to scratch an itch}

\label{lesson:PT-C483-cocar}



\begin{quote}
Before the new one: say the Portuguese for allergic, then the Portuguese for itch.
\end{quote}

\subsection*{You'll want to know: coçar}
\textbf{coçar} — \textquotedblleft{}to scratch an itch\textquotedblright{}.

\begin{quote}
O cão coça a orelha com a pata.
\end{quote}

\begin{quote}
\emph{The dog scratches its ear with its paw.}
\end{quote}

\emph{Coçar-se} — to scratch yourself.

\begin{cousinweb}
\emph{coçar} is inherited, through Old Portuguese, from Vulgar Latin \emph{coctiāre} \textquotedblleft{}to scratch\textquotedblright{} (EtymDB).
\end{cousinweb}

\begin{grammarlens}[title={What you've built}]
One more word for telling, explaining and describing at B1.
\end{grammarlens}

\subsection*{Your turn}
\begin{itemize}
\raggedright
  \item \emph{Say it:} \emph{coçar}
  \item \emph{Say it:} \emph{coçar}, once more
  \item \emph{Say it:} a sentence of your own with \emph{coçar}
  \item \emph{Recall:} say the Portuguese for allergic, then the Portuguese for itch, then say \emph{coçar} again
\end{itemize}

\begin{tcolorbox}[breakable,colback=teal!4,colframe=teal!35!black,title={Before you move on}]
What does coçar mean? (\textquotedblleft{}To scratch an itch\textquotedblright{}.) Say it once more.
\end{tcolorbox}
